%% Chapter 1: Introduction and Problem Statement

\chapter{INTRODUCTION AND PROBLEM STATEMENT}

Unmanned aerial vehicles (UAVs) have emerged as a compelling platform for last-mile
delivery in scenarios where ground transportation is impractical---medical supply chains
to remote clinics, emergency drug drops during disaster response, and fragile diagnostic
sample transport where road vibration risks contamination. In each scenario the UAV
must fly a route from a dispatch hub to a delivery point across an airspace that contains
obstacles, regulated no-fly zones, and variable wind corridors. A good route balances
three conflicting physical objectives simultaneously: minimising total flight distance,
minimising battery energy expenditure, and avoiding high-risk segments. These objectives
are genuinely in tension---the shortest path crosses the most densely occupied airspace,
the safest corridor requires significant detours, and the most energy-efficient altitude
profile rarely coincides with either.

The standard engineering response to such a conflict is to form a weighted sum of the
three objectives into a single scalar cost function, and then run a conventional path
planner such as A* on it. This scalar collapses the multi-objective problem into a
single-objective one, and the operator adjusts the three weights to express the
mission's priority before each flight. The approach works for a fixed mission type,
but it fails immediately when the payload changes: emergency medicine requires a
completely different weight profile from fragile biological samples or disaster-relief
bulk supplies. Manually re-tuning and re-running the planner between missions is
operationally impractical in time-critical field conditions.

This chapter presents the formal problem statement, the research objectives the
proposed system addresses, and the motivation for the two-tier architectural
separation that is the core contribution of this work.

%% ─────────────────────────────────────────────────────────────────────────────
\section{PROBLEM DOMAIN}

Autonomous UAV systems for medical and emergency last-mile delivery operate under a
structural constraint that distinguishes them from general routing problems: the optimal
route trade-off is not fixed---it is a function of what is being carried. A fragile
diagnostic sample demands a low-vibration, low-turbulence corridor even at the cost of
a much longer flight time. Emergency insulin requires the shortest path regardless of
risk. Bulk disaster-relief supplies demand the most energy-efficient route to maximise
delivery range.

No current UAV route planner adjusts its routing priorities to match the payload being
carried. The same fixed objective weights are applied to every mission, forcing field
operators to manually recalibrate between every deployment. This is impractical and
error-prone in time-critical emergency scenarios.

%% ─────────────────────────────────────────────────────────────────────────────
\section{PROBLEM STATEMENT}

Existing UAV route planners encode distance, energy consumption, and flight risk into a
\textbf{single hand-tuned weighted-sum scalar} before path search begins, permanently
fixing the objective trade-off regardless of what is being carried.

When the payload context changes---from fragile diagnostic samples requiring
vibration-minimal routing to time-critical emergency medicine requiring shortest-path
routing---an expert must manually re-tune the planner weights and re-run the search from
scratch. This intervention is \textbf{infeasible in field conditions}.

The proposed system resolves this by separating Pareto front generation from
mission-context-driven route selection. Payload metadata (fragility class, payload
weight class, delivery urgency) has never been used as the sole input to autonomous Pareto
point selection in the UAV delivery literature, leaving a structural gap this project
directly addresses.

%% ─────────────────────────────────────────────────────────────────────────────
\section{MOTIVATION}

The motivation for this work arises from three real operational scenarios that impose
fundamentally incompatible routing requirements:

\textbf{Fragile medical samples} (tissue biopsies, blood samples, pathology specimens)
are highly sensitive to vibration, g-forces, and thermal exposure. The dominant routing
priority is minimising flight risk---defined here as exposure to turbulent corridors,
sudden altitude changes, and wind-exposed segments. A route that adds 40\% extra distance
but halves the risk of vibration damage is operationally preferable.

\textbf{Emergency medicine} (insulin, epinephrine, antivenom) has a strict time window.
The dominant routing priority is minimising route distance to minimise transit time.
Flight risk and energy are secondary. For a payload with a 15-minute delivery window,
the shortest path wins unconditionally.

\textbf{Disaster-relief bulk supplies} (water purification tablets, bandages, ration
packs) weigh significantly more than medical payloads and are not time-critical or
fragile. The dominant routing priority is energy efficiency---minimising battery drain to
maximise the number of deliveries possible on a single charge cycle.

No existing UAV planner can switch between these three routing strategies based on a
payload tag. The proposed two-tier system accomplishes this without any expert
intervention or replanning.

%% ─────────────────────────────────────────────────────────────────────────────
\section{RESEARCH OBJECTIVES}

The following objectives guide this project:

\begin{enumerate}
  \item Formulate multi-objective UAV route planning as a structured \textbf{Pareto front
        pre-computation} problem over a delivery airspace graph, theoretically
        distinguishing the deterministic graph-search approach from stochastic
        population-based alternatives.

  \item Design a \textbf{payload metadata schema} and a principled rule-based mapping
        that translates mission context (fragility class, delivery urgency, payload
        weight class) into Pareto point selection objective weights without any expert
        weight tuning.

  \item \textbf{Benchmark Tier~1} route quality and Pareto front coverage against
        single-objective A* and MOPSO baselines on identical delivery scenarios,
        measuring hypervolume and Pareto spread (Inverted Generational Distance).

  \item \textbf{Quantify Tier~2 selection accuracy} by measuring alignment between
        selected route characteristics and payload tag requirements across diverse
        mission types (fragile medical sample, emergency medicine, disaster-relief
        supplies).

  \item Establish the system's \textbf{novelty}: payload metadata is used as the sole
        input to Pareto point selection, replacing context-blind entropy and
        diversity-based criteria with a deterministic algebraic mapping from payload
        description to route choice.
\end{enumerate}

%% ─────────────────────────────────────────────────────────────────────────────
\section{ORGANISATION OF THE REPORT}

The remainder of this report is organised as follows. Chapter~2 presents a literature
survey of three directly relevant multi-objective UAV path planning papers, identifies
the common research gap they share, and states the research gaps this project addresses.
Chapter~3 describes the proposed two-tier system architecture in detail, including
pseudocode for both tiers, a architecture diagram of the full system, and separate
module descriptions for Phase~I (software) and Phase~II (hardware). Chapter~4 explains
the demonstration plan for both phases. Chapter~5 presents the project timeline as a
Gantt chart. References are listed at the end.